\chapter{LANGUAGE DESIGN AND GRAMMAR}

\section{Types}
% TODO: Describe the four primitive types supported by the language.

\subsection{Integer}
% TODO: Define Integer type, literals, range.

\subsection{Float}
% TODO: Define Float type, literals (e.g. 3.14).

\subsection{Boolean}
% TODO: Define Boolean type, literals (true / false).

\subsection{String}
% TODO: Define String type, literals (e.g. "hello").

\section{Static Typing}
% TODO: Explain static typing: types are checked at compile/parse time,
% not at runtime. No explicit type declaration --- types are inferred.

\section{Token Specification}
% TODO: Table of all tokens with regex patterns:
%   - INT_LITERAL, FLOAT_LITERAL, BOOL_LITERAL, STRING_LITERAL
%   - Identifiers (variable/function names)
%   - Arithmetic operators: +, -, *, /
%   - Comparison operators: ==, !=
%   - Assignment: =
%   - Keywords: if, else, while, def, return, print
%   - Delimiters: (, ), {, }, ,, ;

\section{Context-Free Grammar}
% TODO: Define the formal CFG:
%   G = (V_N, V_T, P, S)
% List all non-terminals and terminals.

\subsection{Program Structure}
% TODO: Top-level production rules (program, statement list, etc.)

\subsection{Arithmetic Expressions}
% TODO: Grammar rules for arithmetic expressions with precedence:
%   expr   -> expr + term | expr - term | term
%   term   -> term * factor | term / factor | factor
%   factor -> INT_LITERAL | FLOAT_LITERAL | IDENTIFIER | ( expr )

\subsection{Boolean Expressions}
% TODO: Grammar rules for boolean expressions (equality and inequality):
%   bool_expr -> expr == expr
%             | expr != expr
%             | BOOL_LITERAL

\subsection{Statements}
% TODO: Grammar rules for:
%   - Assignment:    IDENTIFIER = expr
%   - If-then-else:  if ( bool_expr ) { stmts } else { stmts }
%   - While-loop:    while ( bool_expr ) { stmts }
%   - Print:         print( expr )
%   - Return:        return expr

\subsection{Function Definitions and Calls}
% TODO: Grammar rules for function definition (def keyword, parameter list)
% and function call with value parameter passing.

\section{Operator Precedence and Associativity}
% TODO: Table of operator precedence levels (low to high):
%   Level 1: == , !=   (comparison, non-associative)
%   Level 2: + , -     (addition/subtraction, left)
%   Level 3: * , /     (multiplication/division, left)
% Explain how precedence is encoded in the LALR(1) parser.
